%\chapter{Riadenie a analýza rizík}
%Posúdenie riziká je systematický proces, ktorým presne vyjadrujeme hodnotu rizika (kvantitatívne alebo kvalitatívne),
%t.z. špecifikovať to, čo je ohrozené, vypočítať pravdepodobnosť neželaných zložiek.

%Riziko je veľmi široký pojem, ktorého význam sa v rôznych oblastiach často líši, pre príklad ilustrujeme niekoľko takých príkladov \cite{WikiRisk}.
%V ekonomike znamená riziko neistotu ohľadom straty, vo financiách volatilita výnosu alebo v zdravotníctve možnosť, že nám niečo spôsobí ujmu na zdraví.

%V tejto práci definujeme pojem riziko, ako pravdepodobnosť, že ľudské konanie alebo udalosti vedú k následkom, ktoré majú vplyv na to, čo si ľudia cenia \cite{Ras83}.


% Na začiatok by sme čítateľa radi upozornili na nekonzistentnosť termínov naprieč rôznymi štandardami, čo môže vyústiť v zbytočný zmätok a chaos. V tejto práci sa
\chapter{Analýza rizík}
\section{Analýza rizík podľa BSI 200-3}
Analýza rizík je proces, ktorým ustanovujeme a ošetrujeme riziko. Podľa \cite{BSI3} tento proces pozostáva zo štyroch hlavných krokov a
niekoľkých podkrokov (viď \ref{fig:risk-analysis-bsi3}).
V časti \ref{sec:bsiphilosophy} sme naznačili, že BSI sa snaží \say{odbremeniť} malé a stredne veľké organizácie, čo môžeme vidieť aj v
v tom, ako postupovali pri analýze rizík: rozdelili ju na dve časti, a to implicitnú a explicitnú.

\input tikz/risk-analysis-bsi3.tex

\textbf{Implicitná} časť analýzy rizík je tá, ktorú za nás vykonal BSI a je súčasťou IT-Grundshutz kompendia \cite{GSK}. BSI vybral
najčastejšie\footnote{Podľa \cite{mngmtkib} je to 80\% spoločných aktív, zvyšných 20\% tvoria aktíva špecifické pre konkrétnu organizáciu.}
aktíva, ktoré používa väčšina organizácii a pre ne vykonali analýzu rizík. Všetky potrebné informácie (popis, cieľ, rozsah, opatrenia) sú
detailne obsiahnuté v moduloch \cite{GSK}.

\textbf{Explicitná} časť analýzy rizík pokrýva tie aktíva, ktoré nemôžeme namapovať na žiaden modul \cite{GSK} alebo bezpečnostné opatrenia pre
namapovaný modul \cite{GSK} neposkytujú dostatočnú ochranu. Návod, ako môžeme vykonať explicitnú analýzu rizík pre aktívum popisuje \cite{BSI3}.
Vo zvyšnej časti práce budeme pojmom analýza rizík referovať na \emph{explictnú analýzu rizík}, pokiaľ explicitne neuvedieme, že referejueme na
\emph{implicitnú analýzu rizík}.

V minulej časti (\ref{sec:isms}) sme popísali kroky vedúce k analýze rizík, ktoré za pomoci \cite{BSI2} vieme zhrnúť na:
\begin{itemize}[nosep]
    \item inicializácia systematického procesu KIB
    \item určenie rozsahu bezpečnostného konceptu, t.j. určenie informačného systému
    \item vypracovanie štrukturovanej analýzy informačného systému
    \item posúdenie bezpečnostných požiadaviek na aktíva
    \item modelovanie aktív na moduly IT-Grundschutz \cite{GSK}
    \item vykonanie kontroly IT-Grundschutz
\end{itemize}

Postupnou implementáciou bezpečnostného konceptu sme sa dopracovali až ku kroku \emph{IT-Grundschutz check}, z ktorého nám vypadli tzv. cieľové objekty
podliehajúce kontrole (\emph{angl. target objects under review}). Inými slovami sú to cieľové aktíva, voči ktorým potrebujeme explicitne vykonať
analýzu rizík. Teraz prejdeme k detailom tohto procesu a vysvetlíme podrobne jednotlivé časti obrázku \ref{fig:risk-analysis-bsi3}, pre lepšie vysvetlenie 
si zadefinujeme modelovú situáciu. Upozorňujeme autora, že realitu a detaily mierne skreslíme pre jednoduchosť ilustrácie. Uvedomujeme si potenciálnu
náročnosť a komplexnosť problémov

\subsection*{Fiktívny scenár}

Predpokladajme organizáciu Národný olympijský inštitút pre pečenie (NOIP), ktorého primárnou agendou je zabezpečovať organizáciu rôznych kôl a kategórií
pečenia naprieč Slovenskom. Nakoľko NOIP je štátna inštitúcia, vzťahujú sa na ňu relevantné zákony KIB (viď kapitolu \ref{chap:zakony}), a preto za pomoci 
manažéra KIB sme iniciovali bezpečnostný proces.

NOIP je moderná organizácia a jej zamestnanci sú vybavení pracovnými notebookmi, pomocou ktorých vykonávajú svoju pracovnú náplň, ako napríklad komunikácia s kolegami,
organizátormi a súťažiacimi. Počas analýzy procesov a bezpečnostných požiadaviek sme dospeli k tomu, že notebooky sú kľúčovým aktívom pre zamestnancov. Špecifická dynamika
organizovania olympiád (cestovanie a prítomnosť na podujatiach) a administratíva (spracovanie osobných údajov a utajenie úloh) nás prinútila toto aktívum
podrobiť explicitnej analýze rizík. Notebooky sme vo fáze modelovania namapovali na modul \cite{GSK} \emph{SYS.3.1 Laptops}.

\subsection{Určenie relevantných hrozieb}
Pri určovaní relevantných hrozieb vychádzame z modulu, na ktorý sme aktívum namapovali. Postupujeme v troch krokoch, najprv z katalógu \cite{GSK} pre relevantný modul
zozbierame identifikované základné hrozby, potom posúdime relevantnosť základných hrozieb z katalógu a identifikujeme dodatočné špecifické hrozby.

Relevantnosť základnej hrozby vyjadrujeme v troch mierach irelevantná, nepriamo relevantná a relevantná (pozri vysvetlenie \ref{tab:elemthreatrel}).
\begin{table}
    \caption{Relevantnosť základnej hrozby}
    \label{tab:elemthreatrel}
    \begin{center}
        \begin{tabular}{ |c|c|c| } 
         \hline
         \textbf{druh hrozby} & \textbf{dopad na aktívum} & \textbf{ďalší postup} \\ 
         \hline
         \hline
         \textbf{irelevantná} & nemá žiaden & vylúčená \\ 
         \hline
         \textbf{nepriamo relevantná} & cez iné aktívum & vylúčená, ošetrená priamo iným opatrením  \\ 
         \hline
         \textbf{relevantná} & priamy & zaradená do analýzy rizík \\ 
         \hline
        \end{tabular}
    \end{center}
\end{table}
Dodatočné špecifické hrozby vyplývajú z unikátnych podmienok, v ktorých sa aktívum nachádza. Pred tým, ako sformulujeme špecifickú hrozbu,
by sme mali zvážiť to, či nová hrozba má potenciál spôsobiť významnú škodu a či je realistická v kontexte využitia nášho aktíva\footnote{
Množinu relevantných hrozieb sa snažíme minimalizovať, aby sme udržali rozsah na rozumnej a zvládnuteľnej úrovni s ohľadom na ďalšie kroky
analýzy rizík.}.

\subsection*{Určenie relevantných hrozieb: fiktívny scenár}

Pri podrobnejšej inšpekcii procesu olympiád sme zistili, že pri výjazde na súťažné kolo si zamestnanci nesú pracovné notebooky so sebou.
Počas tohto výjazdu prichádzajú do styku s mnohými účastníkmi sútaže a technickými prvkami, ktoré majú potenciál negatívneho dopadu na
základné hodnoty CIA notebooku. Určenie relevantných základných hrozieb sme zhrnuli\footnote{Hrozby už identifikované \cite{GSK}, ani
irelevantné hrozby explicitne neuvádzame.} do tabuľky \ref{tab:elemthreatid}.
\begin{table}[h!]
    \caption{Identifikované základné hrozby}
    \label{tab:elemthreatid}
    \begin{center}
        \footnotesize
        \begin{tabular}{ |p{0.35\textwidth}|p{0.09\textwidth}|p{0.56\textwidth}| }
        \hline
        hrozba & dopad & poznámka \\
        \hline
        \hline
        G 0.25 Sociálne inžinierstvo & nepriamy & Hrozba G 0.25 by umožnila naplniť hrozby G 0.19 Zverejnenie citlivých informácií alebo G 0.30,
        no samotný laptop neposkytuje žiadne bezpečnostné opatrenia pred sociálnym inžinierstvom, preto ju považujeme za nepriamu.  \\
        \hline
        G 0.30 Neoprávnené používanie alebo správa zariadení a systémov & priamy &  Hrozba G 0.30 priamo vplýva na laptop, preto musíme
        skontrolovať bezpečnostné opatrenia voči G 0.30.\\
        \hline
        \end{tabular}
    \end{center}
\end{table}

Podobne sme postupovali pri identifikácií dodatočných špecifických hrozieb v tabuľke \ref{tab:specthreatid}. V popise explicitne uvádzame,
z ktorých ďalších hrozieb vychádzame, čo ale nie je potrebné.
\begin{table}[h!]
    \caption{Identifikované špecifické hrozby}
    \label{tab:specthreatid}
    \begin{center}
        \footnotesize
        \begin{tabular}{ |p{0.20\textwidth}|p{0.80\textwidth}| }
        \hline
        hrozba & popis \\
        \hline
        \hline
        T z.1 Poškodenie zariadení v súťažnom prostredí &
        \parbox[t]{0.80\textwidth}{
        Notebooky sa vyskytujú v prítomnosti chaotického súťažného prostredia, kde sú vystavené náhlym a rýchlym pohybom alebo koncentrovanému
        množstvu tekutín a prachových častíc. \\\\ Táto dodatočná špecifická hrozba rozvíja základnú hrozbu G 0.4 Znečistenie, prach, korózia 
        }\\

        \hline
        T z.2 Manipulácia v súťažnom prostredí &
        \parbox[t]{0.80\textwidth}{
        V súťažnom prostredí môžu účastníci (učitelia, súťažiaci, pomocníci) dočasne získať príležitosť nedovoleného prístupu k notebooku, čím
        môžu kompromitovať integritu alebo dostupnosť dôverných informácií. \\\\ Táto dodatočná špecifická hrozba rozvíja základne hrozby 
        G 0.19 Zverejnenie citlivých informácií a G 0.22 Manipulácia s informáciami 
        }\\
        \hline
        \end{tabular}
    \end{center}
\end{table}
Keď máme zostavený zoznam relevantných hrozieb, ďalším krokom je klasifikácia rizík, ktoré z identifikovaných hrozieb vyplývajú.

\section{Klasifikácia rizík}

Najčastejšie vyjadrujeme riziko ako strednú hodnotu dopadu hrozby, t.j. $riziko = pravdepodobnosť * dopad hrozby$ \cite{MFSRIB}, neexistuje
univerzálny spôsob, ako určiť hodnotu týchto premenných. Pri určovaní dopadu môžeme brať do úvahy potenciálne finančné následky, priamy a
latentný dopad na okolie alebo úsilie potrebne na nápravu a odstránenie škody, pri určovaní pravdepodobnosti môžem vychádzať z vlastných skúseností 
alebo štatistiky\footnote{Treba si však dať pozor, nakoľko výsledky štatistiky sú závisle od kontextu a interpretácia jej výsledkov  môže
byť subjektívna.}. Vo všeobecnosti na vyjadrenie rizika používame dve metódy kvalitatívnu a kvantitatívnu, štandardy BSI využívajú tú prvú,
no vysvetlíme oboje.

\textbf{Kvantitatívna metóda} vyjadruje hodnotu číselne ($interval <0,1>$), avšak problémom je presne číselne vyjadriť pravdepodobnosť a dopad
v špecifickom kontexte našej organizácie, problematické sa môžu javiť nasledovné otázky \cite{MFSRIB}:
\begin{itemize}[nosep]
    \item ako vyjadriť pravdepodobnosť udalosti, ktorá ešte nikdy nenastala?
    \item ako vyjadriť hodnotu nehmotných aktív?
\end{itemize}

\textbf{Kvalitatívna metóda} vyjadruje hodnotu slovne, IT-Grundschutz používa $4x4x4$ úrovní, v poradí pravdepodobnosť (\ref{tab:thrfreq}),
dopadu (\ref{tab:thrdmg}) a riziko (\ref{tab:thrimp}). Prirodzene, v organizácii si počet a význam úrovní potrebujem prispôsobiť svojím vlastným
podmienkam a potrebám. 
\begin{table}[!h]
    \caption{Frekvencia výskytu}
    \label{tab:thrfreq}
    \begin{center}
        \footnotesize
        \begin{tabularx}{\textwidth}{|X|X|}
        \hline
        úroveň & popis \\
        \hline
        \hline
        zriedkavo & najviac každých 5 rokov \\
        \hline
        stredne & raz za 1-5 rokov \\
        \hline
        často & raz za mesiac \\
        \hline
        veľmi často & niekoľkokrát do mesiaca  \\
        \hline
        \end{tabularx}
    \end{center}
\end{table}

\begin{table}[!h]
    \caption{Potenciálny dopad}
    \label{tab:thrdmg}
    \begin{center}
        \footnotesize
        \begin{tabularx}{\textwidth}{|X|X|}
        \hline
        úroveň & popis \\
        \hline
        \hline
        zanedbateľný & účinky poškodenia sú nízke a môžme ich zanedbať  \\
        \hline
        limitovaný & účinky poškodenia sú obmedzené ale zvládnuteľné \\
        \hline
        signifikantný & účinky poškodenia sú významné  \\
        \hline
        života-ohrozujúci & účinky poškodenia môžu byť katastrofálne a ohrozovať existenciu organizácie \\
        \hline
        \end{tabularx}
    \end{center}
\end{table}

\begin{table}[!h]
    \caption{Riziko}
    \label{tab:thrimp}
    \begin{center}
        \footnotesize
        \begin{tabularx}{\textwidth}{|X|X|}
        \hline
        úroveň & popis \\
        \hline
        \hline
        nízke & navrhnuté alebo implementované bezpečnostné opatrenia poskytujú dostatočnú ochranu \\
        \hline
        stredné & navrhnuté alebo implementované bezpečnostné opatrenia nemusia byť dostatočné \\
        \hline
        vysoké & navrhnuté alebo implementované bezpečnostné opatrenia poskytujú dostatočnú ochranu \\
        \hline
        veľmi vysoké & navrhnuté alebo implementované bezpečnostné opatrenia poskytujú dostatočnú ochranu, v praxi sú málokedy akceptované \\
        \hline
        \end{tabularx}
    \end{center}
\end{table}

Nakoniec vzťah pravdepodobnosti, potenciálnej škody a potenciálneho dopadu vyjadríme v maticou klasifikácie rizík \ref{fig:risk-matrix-bsi3}.
\definecolor{cellgreen}{RGB}{34,170,79}
\definecolor{celllightgreen}{RGB}{141,198,63}
\definecolor{cellyellow}{RGB}{255,237,0}
\definecolor{cellred}{RGB}{228,30,38}

\begin{figure}[h!]
    \centering
    \begin{tikzpicture}

        % Cell size (halved from original 3.0/2.0)
        \def\cw{1.5}
        \def\ch{1.5}
        % Padding between cells
        \def\pad{0.08}

        % Helper: draw a padded filled rectangle
        % #1=color, #2=col(0-3), #3=row(0-3)
        \newcommand{\cell}[3]{%
            \fill[#1, rounded corners=1pt]
            ({#2*\cw+\pad}, {#3*\ch+\pad})
            rectangle
            ({(#2+1)*\cw-\pad}, {(#3+1)*\ch-\pad});
            \draw[black, rounded corners=1pt]
            ({#2*\cw+\pad}, {#3*\ch+\pad})
            rectangle
            ({(#2+1)*\cw-\pad}, {(#3+1)*\ch-\pad});
        }

        % --- Row 1 (bottom): negligible ---
        \cell{cellgreen}{0}{0}
        \cell{cellgreen}{1}{0}
        \cell{cellgreen}{2}{0}
        \cell{cellgreen}{3}{0}

        % --- Row 2: limited ---
        \cell{cellgreen}{0}{1}
        \cell{cellgreen}{1}{1}
        \cell{celllightgreen}{2}{1}
        \cell{cellyellow}{3}{1}

        % --- Row 3: significant ---
        \cell{celllightgreen}{0}{2}
        \cell{celllightgreen}{1}{2}
        \cell{cellyellow}{2}{2}
        \cell{cellred}{3}{2}

        % --- Row 4 (top): life-threatening ---
        \cell{celllightgreen}{0}{3}
        \cell{cellyellow}{1}{3}
        \cell{cellred}{2}{3}
        \cell{cellred}{3}{3}

        % --- Cell labels ---
        % Row 1
        \node[font=\scriptsize\bfseries, text=white] at (0.5*\cw, 0.5*\ch) {nízke};
        \node[font=\scriptsize\bfseries, text=white] at (1.5*\cw, 0.5*\ch) {nízke};
        \node[font=\scriptsize\bfseries, text=white] at (2.5*\cw, 0.5*\ch) {nízke};
        \node[font=\scriptsize\bfseries, text=white] at (3.5*\cw, 0.5*\ch) {nízke};
        % Row 2
        \node[font=\scriptsize\bfseries, text=white] at (0.5*\cw, 1.5*\ch) {nízke};
        \node[font=\scriptsize\bfseries, text=white] at (1.5*\cw, 1.5*\ch) {nízke};
        \node[font=\scriptsize\bfseries] at (2.5*\cw, 1.5*\ch) {stredné};
        \node[font=\scriptsize\bfseries] at (3.5*\cw, 1.5*\ch) {vysoké};
        % Row 3
        \node[font=\scriptsize\bfseries] at (0.5*\cw, 2.5*\ch) {stredné};
        \node[font=\scriptsize\bfseries] at (1.5*\cw, 2.5*\ch) {stredné};
        \node[font=\scriptsize\bfseries] at (2.5*\cw, 2.5*\ch) {vysoké};
        \node[font=\scriptsize\bfseries, text=white, align=center] at (3.5*\cw, 2.5*\ch) {veľmi\\vysoké};
        % Row 4
        \node[font=\scriptsize\bfseries] at (0.5*\cw, 3.5*\ch) {stredné};
        \node[font=\scriptsize\bfseries] at (1.5*\cw, 3.5*\ch) {vysoké};
        \node[font=\scriptsize\bfseries, text=white, align=center] at (2.5*\cw, 3.5*\ch) {veľmi\\vysoké};
        \node[font=\scriptsize\bfseries, text=white, align=center] at (3.5*\cw, 3.5*\ch) {veľmi\\vysoké};

        % --- Y-axis ---
        \draw[thick, -stealth] (-0.1, -0.1) -- (-0.1, 4*\ch+0.4);
        \node[anchor=south, font=\scriptsize\bfseries, align=center] at (-0.1, 4*\ch+0.4) {potenciálny dopad};

        \node[anchor=east, font=\scriptsize] at (-0.2, 0.5*\ch) {zanedbateľný};
        \node[anchor=east, font=\scriptsize] at (-0.2, 1.5*\ch) {limitovaný};
        \node[anchor=east, font=\scriptsize] at (-0.2, 2.5*\ch) {signifikantný};
        \node[anchor=east, font=\scriptsize, align=right] at (-0.2, 3.5*\ch) {života-\\ohrozujúci};

        % --- X-axis ---
        \draw[thick, -stealth] (-0.1, -0.1) -- (4*\cw+0.4, -0.1);
        \node[anchor=west, font=\scriptsize\bfseries, align=center] at (4*\cw+0.4, -0.2) {frekvencia\\výskytu};

        \node[anchor=north, font=\scriptsize] at (0.5*\cw, -0.15) {zriedkavo};
        \node[anchor=north, font=\scriptsize] at (1.5*\cw, -0.15) {stredne};
        \node[anchor=north, font=\scriptsize] at (2.5*\cw, -0.15) {často};
        \node[anchor=north, font=\scriptsize] at (3.5*\cw, -0.15) {veľmi často};

    \end{tikzpicture}
    \caption{Matica rizík, na osi x frekvencia výskytu (rastie v smere šípky), na osi y potenciálny dopad (rastie v smere šípky), bunka vyjadruje potenciálny dopad hrozby \cite{BSI3}.}
    \label{fig:risk-matrix-bsi3}
\end{figure}

\subsection*{Klasifikácia rizík: fiktívny scenár}
Vychádajúc z fázy bezpečnostného procesu, kde sme definovali bezpečnostné požiadavky aktív, sme vytvorili tabuľku popisu hrozieb. Z predchádzajúcich
skúsenosti sme odhadli frekvenciu výskytu a potenciálny dopad a v kontexte načrtnutých scenárov sme objasnili vyplývajúce riziko v tabuľke
\ref{tbl:riskassesmentexample}, pre prehľadnosť sme uviedli iba jednu klasifikáciu hrozby. Upozorníme na to, že nie vždy sú potrebné rozsiahle popisy
a vysvetlenia, najmä ak sú vo svojej podstate triviálne. Písanie detailných opisov a zdôvodnení sa môže stať časovo netriviálnou záležitosťou v prípade
rozsiahleho počtu aktív a hrozieb.

\begin{table}[h!]
\label{tbl:riskassesmentexample}
\definecolor{headerblue}{RGB}{220, 225, 235}
\definecolor{riskyellow}{RGB}{255, 192, 0}

\begin{tabularx}{\textwidth}{|X|}
\hline
\rowcolor{headerblue}
\textbf{Notebooky N1} \\
\rowcolor{headerblue}
Dôvernosť: veľmi vysoká \\
\rowcolor{headerblue}
Integrita: vysoká \\
\rowcolor{headerblue}
Dostupnosť: veľmi vysoká \\
\hline
\end{tabularx}

\begin{tabularx}{\textwidth}{|X|X|}
\hline
Hrozba \newline
G 0.30 \textit{Neoprávnené používanie alebo správa zariadení a systémov}
&
Narušené základné hodnoty: \newline
dôvernosť, integrita, dostupnosť \\
\hline
\end{tabularx}

\begin{tabularx}{\textwidth}{|X|X|X|}
\hline
Frekvencia výskytu bez dodatočných bezpečnostných opatrení: často
&
Potenciálny dopad bez dodatočných bezpečnostných opatrení: značné
&
\rowcolor{riskyellow}
Riziko bez dodatočných bezpečnostných opatrení: vysoké \\
\hline
\end{tabularx}

\begin{tabularx}{\textwidth}{|X|}
\hline
\textbf{Popis:} \newline
Zamestnanci si pravidelne nosia notebooky na súťažné kolá, kde riadia registráciu a pomáhajú organizátorom podujatia a
z času na čas nakrátko odbehnú. Na notebookoch majú uložené rôzne údaje a materiály, častokrát aj správne súťažné
riešenia pre súťažne kolo, v prípade potreby tlače. Navyše, cez webový prehliadač pristupujú k intranetu organizácie,
kde sú vystavené mnohé neverejné dokumenty.
\newline\newline
\textbf{Vyhodnotenie:} \newline
Administrátori informačných systémov musia zaviesť povinné uzamykanie obrazovky, a to aj zavedením časového limitu pre automatické uzamknutie
zariadenia. Dodatočne by mali zaviesť šifrovanie diskov, aby predišli strate dát a úniku citlivých informácií v prípade narušenia integrity
zariadenia.
\hline
\end{tabularx}
\end{table}
\newpage

\section{Ošetrenie rizík}

Rizika vyplývajúce z hrozieb potrebujeme znížiť na akceptovateľnú úroveň, ktorú sme si určili v politike informačnej bezpečnosti. V našom prípade
predpokladáme akceptovateľnú úroveň rizika \enquote{nízka}, tým pádom nás zaujímajú \enquote{stredné} a vyššie riziká. K dispozícií máme štyri
rozličné mechanizmy, ako znížiť riziko: vyhnúť sa, redukovať, preniesť a akceptovať. Otázky v tabuľke \ref{tab:trmqst} nám môžu pomôcť pri výbere
druhu ošetrenia. 

\textbf{Vyhnúť sa} riziku znamená odstrániť príčinu rizika alebo nositeľa rizika. Tento prístup si môžme zvoliť v prípade, že zavedenie efektívnych
protiopatrení je príliš nákladné alebo značne obmedzujú prístup/správu systému. Tiež môžeme zvážiť reštrukturalizovať procesy a znížiť tým náklady.

\textbf{Redukovať} riziko znamená vyvinúť a zaviesť efektívne opatrenia zničujúce pravdepodobnosť alebo dopady hrozby. Opatrenia môžeme vyvíjať
na základe predchádzajúcich skúsenosti, štandardov, \emph{best-practices} alebo za podpory odborníkov a špecialistov v odvetví.

\textbf{Preniesť} riziko vieme na tretiu stranu istou formou poistenia alebo \emph{outsourcingu}. V zásade ide o finančné škody, kapacitné nedostatky
alebo nedostatočný/chýbajúci \emph{know-how} na zavedenie vlastného riešenia. Pri tejto forme ošetrenia rizika si potrebujeme vhodne nastaviť zmluvné
podmienky alebo zmeniť cez dodatky v prípade už uzavretých zmlúv.


\textbf{Akceptovať} riziko znamená stotožniť sa s potenciálnymi následkami, ktoré následne monitorujeme ako zostatkové riziko. V niektorých prípadoch
je riziko prirodzenou súčasťou náplne organizácie a jeho akceptovanie je súčasťou stratégie. Podobne existujú hrozby, voči ktorým neexistujú efektívne
protiopatrenia alebo náklady na jej minimalizovanie presahujú hodnotu aktíva.



\begin{table}[h!]
    \caption{Otázky, ako pomôcka pri výbere druhu ošetrenia}
    \label{tab:trmqst}
    \begin{center}
        \footnotesize
        \begin{tabularx}{\textwidth}{|l|X|}
        \hline
        vyhnúť sa & Má zmysel vyhnúť sa riziku reštrukturalizovaním biznisového procesu alebo informačného systému?
        \\
        \hline
        redukovať & Má zmysel a je možné znížiť riziko zavedením dodatočných bezpečnostných opatrení?
        \\
        \hline
        preniesť & Má zmysel preniesť riziko na inú organizáciu pomocou poistnej zmluvy alebo outsourcingom?
        \\
        \hline
        akceptovať & Môže byť riziko akceptované na základe zrozumiteľných faktov?
        \\
        \hline
        \end{tabularx}
    \end{center}
\end{table}

\section{Ošetrenie rizík: fiktívny scenár}

Počas brainstormingu sme uvažovali nad vyhnutím sa rizikám laptopov a zákazom ich nosenia na súťažné kolá, čo ale viedlo
k značnému odporu zamestnancov. Preto vzhľadom na nenáročnosť a jednoduchosť implementácie sme sa rozhodli redukovať
riziko dodatočnými opatreniami zhrnutými v tabuľke \ref{tbl:risktreatmentexample}.

\begin{table}[h!]
\definecolor{headerblue}{RGB}{220, 225, 235}
\definecolor{riskyellow}{RGB}{255, 192, 0}
\definecolor{riskgreen}{RGB}{20, 220, 25}

\begin{tabular}{|p{\textwidth}|}
\hline
\rowcolor{headerblue}
\textbf{Notebooky N1} \\
\rowcolor{headerblue}
Dôvernosť: veľmi vysoká \\
\rowcolor{headerblue}
Integrita: vysoká \\
\rowcolor{headerblue}
Dostupnosť: veľmi vysoká \\
\hline
\end{tabular}

\begin{tabular}{|p{0.20\textwidth}|p{0.20\textwidth}|p{0.60\textwidth}|}
\hline
Hrozba & Riziko & Ošetrenie
\\
\hline
G 0.30 \textit{Neoprávnené používanie alebo správa zariadení a systémov}
&
\colorbox{riskyellow}{vysoké}
\newline\newline
S dodatočnými opatreniami:
\newline
\colorbox{riskgreen}{nízke}
& redukcia rizika
\newline
Doplňujúce bezpečnostné opatrenia:
Notebooky majú zavedené zamknutie obrazovky a vypnutie po uplynutí časového limitu. Dodatočne, notebooky majú
šifrované disky a vyžadujú si hardvérový kľuč pri jeho zapínaní.
\\
\hline
\end{tabular}
\caption{Ošetrenie rizika vo fiktívnej organizácii}
\label{tbl:risktreatmentexample}
\end{table}

\newpage

\section{Konsolidácia}

V prípade, že zavedieme dodatočné bezpečnostné opatrenia, musíme konsolidovať\footnote{Slovníková definícia, uviesť do
poriadku (pomery, vzťahy a pod.), stabilizovať, upevniť, ustáliť.} bezpečnostný koncept. Potrebujeme posúdiť, či nové
bezpečnostné opatrenia sú vhodné na neutralizáciu hrozby, ako ovplyvňujú už zavedené bezpečnostné opatrenia, či sú
užívateľsky prívetivé a jasné, alebo či sú vynaložené prostriedky vhodne použité. Nevhodné a neefektívne bezpečnostné
opatrenia vyradíme alebo nahradíme novými, nepostačujúce doplníme a vylepšíme.

\section{Definovanie bezpečnostných požiadaviek}

Vhodne definované bezpečnostné požiadavky voči základným hodnotám CIA aktíva identifikovaného v informačnej doméne tvoria podstatnú časť
bezpečnostného procesu, kedy si vysoko-ohrozené aktíva vyžadujú včasnú pozornosť a prednosť, čo môže v prvom rade ovplyvniť priority manažéra
informačnej bezpečnosti alebo aj uprednostniť voľbu jadrovej ochrany pred základnou alebo štandardnou ochranou, v neposlednom rade aj počet aktív,
vstupujúcich do explicitnej analýzy rizík. Na základe týchto dôvodov sme sa rozhodli venovať časť tejto práce aj tomu, ako (ne)vhodne definovať
bezpečnostné požiadavky pre dané aktívum.

Bezpečnostné požiadavky na aktívum sú zamerané na potenciálne škody, ktoré vyplývajú z nepriaznivého vplyvu objektívne existujúcej hrozby. Podľa
\cite{BSI3} uvažujeme tri kategórie ochrany:
\begin{itemize}[nosep]
    \item normálna, dôsledky narušenia sú limitované zvládnuteľne 
    \item vysoká, dôsledky narušenia sú značné
    \item veľmi vysoká, dôsledky narušenia môžu mať katastrofálne dopady a ohrozovať existenciu organizácie 
\end{itemize}
Počet kategórií nie je vopred daný, môže sa líšiť aj naprieč CIA a musíme si ich prispôsobiť vlastným podmienkam a potrebám\footnote{Napr. môžeme
dôvernosť rozdeliť na štyri úrovne: verejné, interné, dôverné a utajené; dostupnosť na dve úrovne: normálna a kritická.}. Zoberme kategóriu
biznisový proces, kde narušenia CIA vieme rozdeliť do niekoľkých spoločných scenárov:
\begin{itemize}[nosep]
    \item porušenie zákonov, regulácií a zmluvných podmienok
    \item porušenie práva na informačne sebaurčenie \footnote{\emph{Recht auf informationelle Selbstbestimmung},
    právo jednotlivca, aby na základe princípu sebaurčenia sám rozhodoval, kedy a v akom rozsahu by mali byť
    informácie o jeho súkromnom živote poskytované iným.}
    \item narušenie fyzickej integrity osoby
    \item narušenie schopnosti vykonávať úlohu
    \item negatívne interné alebo externé účinky
    \item finančné dôsledky
\end{itemize}
Potom pri posudzovaní úrovne ochrany môžme postupovať podľa kritérií v tabuľke \ref{tab:protectionscenario}, neuvádzali sme všetky skupiny, vybrali sme
tie najzaujímavejšie príklady.


\begin{table}[h!]
    \caption{Vybrané scenáre na ilustráciu posudzovania úrovne ochrany}
    \label{tab:protectionscenario}
    \begin{center}
            \begin{tabularx}{\textwidth}{|>{\raggedright\arraybackslash}p{2.8cm}
                                  |>{\raggedright\arraybackslash}X
                                  |>{\raggedright\arraybackslash}X
                                  |>{\raggedright\arraybackslash}X|}
            \hline
            \textbf{scenár} & \textbf{normálna ochrana} & \textbf{vysoká ochrana} & \textbf{veľmi vysoká ochrana}
            \\
            \hline
            \hline
            \textbf{porušenie práva na informačne sebaurčenie}
             &
            narušenie tých osobných údajov, ktorých spracovanie môže mať nepriaznivé účinky na spoločenské postavenie alebo ekonomické podmienky dotknutej osoby
             &
            narušenie tých osobných údajov, ktorých spracovanie môže mať značne nepriaznivé účinky na spoločenské postavenie alebo ekonomické podmienky dotknutej osoby
             &
            ohrozenie života a zdravia alebo osobných slobôd dotyčnej osoby
            \\
            \hline
            \textbf{narušenie schopnosti vykonávať úlohu}
             &
            narušenie by bolo posúdené ako znesiteľné, maximálny akceptovateľný prestoj je medzi 24 až 72 hodín
             &
            narušenie by bolo posúdené ako neznesiteľné, maximálny akceptovateľný prestoj je medzi hodinou až 24 hodín
             &
            narušenie by bolo posúdené ako neakceptovateľné, maximálny akceptovateľný prestoj je menej, ako jedna hodina
            \\
            \hline
            \textbf{finančné dôsledky}
             &
            akceptovateľná finančná strata
             &
            značná finančná strata, ktorá neohrozuje existenciu organizácie
             &
            finančná strata, ktorá ohrozuje existenciu organizácie
            \\
            \hline
        \end{tabularx}
    \end{center}
\end{table}

Poslednú praktickú radu, ktorú uvedieme, je sada štyroch princípov, ktorými by sme sa mali riadiť pri určovaní
bezpečnostných požiadaviek:
\begin{itemize}[nosep]
    \item \textbf{dedenie}, postupujeme smerom z vrchu, biznisové procesy a aplikácie, potom ich individuálne prvky, ako miestností, informačné systémy a komunikačné linky
    \item \textbf{maximálnosti}, ak rozdelíme aktívum na menšie častí, berieme do úvahy ten najzávažnejší účinok
    \item \textbf{závislosti}, narušenie podporného aktíva môže spôsobiť narušenie k nemu primárneho aktíva
    \item \textbf{kumulatívne účinky}, veľké množstvo malých narušení môže v konečnom dôsledku vyústiť do jedného veľkého narušenia iného aktíva
    \item \textbf{distribučný efekt}, narušenie primárneho aktíva sa nemusí plne preniesť na jeho sekundárne aktíva
\end{itemize}

V prípade hlbšieho záujmu čitateľa ponúka štandard \cite{BSI2} v kapitole 8.2 kompletné príklady a úvahy pri určovaní
bezpečnostných požiadaviek pre biznisové procesy, aplikácie, IT systémy, miestnosti, komunikačné linky a
iné/nezaradené zariadenia.